\textbf{Persistent Segment Tree}
\begin{lstlisting}
struct PST {
  struct Node { int l, r, val; };
  vector<Node> t;
  int T = 0;
  PST(int max_nodes) { // estimate 20 * n for
  // max_nodes
    t.resize(max_nodes);
  }
  int build(int b, int e) {
    int cur = T++;
    if (b == e) return cur;
    int m = (b + e) / 2;
    t[cur].l = build(b, m);
    t[cur].r = build(m + 1, e);
    return cur;
  }
  int update(int pre, int b, int e, int i, int v) {
    int cur = T++;
    t[cur] = t[pre];
    if (b == e) {
      t[cur].val += v;
      return cur;
    }
    int m = (b + e) / 2;
    if (i <= m) t[cur].l = update(t[pre].l, b, m, i,
    v);
    else t[cur].r = update(t[pre].r, m + 1, e, i, v);
    t[cur].val = t[t[cur].l].val + t[t[cur].r].val;
    return cur;
  }
  int query(int pre, int cur, int b, int e, int k) {
    if (b == e) return b;
    int cnt = t[t[cur].l].val - t[t[pre].l].val;
    int m = (b + e) / 2;
    if (cnt >= k) return query(t[pre].l, t[cur].l, b,
    m, k);
    return query(t[pre].r, t[cur].r, m + 1, e, k - cnt);
  }
};
int main() {
  ios::sync_with_stdio(0); cin.tie(0);
  int n, q;
  cin >> n >> q;
  vector<int> a(n + 1), root(n + 1), V;
  map<int, int> comp;
  for (int i = 1; i <= n; i++) cin >> a[i],
  comp[a[i]];
  int id = 0;
  V.resize(comp.size() + 1);
  for (auto& [x, _] : comp) comp[x] = ++id, V[id]
  = x;
  PST pst(20 * (n + 1)); // Allocate enough space
  root[0] = pst.build(1, id);
  for (int i = 1; i <= n; i++)
  root[i] = pst.update(root[i - 1], 1, id,
  comp[a[i]], 1);
  while (q--) {
    int l, r, k;
    cin >> l >> r >> k;
    cout << V[pst.query(root[l - 1], root[r], 1, id,
    k)] << '\n';
  }
  return 0;
}
\end{lstlisting}
%the code returns k-th number in a range l to r if
%the range were sorted
\begin{lstlisting}
int V[N], root[N], a[N];
int32_t main() {
  map<int, int>mp;
  int n, q;
  cin >> n >> q;
  for(int i = 1; i <= n; i++) cin >> a[i], mp[a[i]];
  int c = 0;
  for(auto x : mp) mp[x.first] = ++c, V[c] = x.first;
  root[0] = t.build(1, n);
  for(int i = 1; i <= n; i++) {
    root[i] = t.upd(root[i - 1], 1, n, mp[a[i]], 1);
  }
  while(q--) {
    int l, r, k;
    cin >> l >> r >> k;
    cout << V[t.query(root[l - 1], root[r], 1, n, k)]
    << '\n';
  }
  return 0;
}
\end{lstlisting}
